\subsection{I-37. Two centered pairs for affine singular width-three words}
\label{I-37}

\paragraph{Statement and exact promise.}
Let $E$ be any field, let $a_1,\ldots,a_n\in E$ be distinct, and let
$\lambda_i(Y)$ be their Lagrange polynomials. A nonzero polynomial
$f\in E[X_1,\ldots,X_n]$ is assumed to admit a width-at-most-three
ROABP whose variable-dependent factors are affine-linear in their
own variables and have generic rank at most two. Constant factors
are allowed. The order is arbitrary and unknown to the map. Then
\[
 f\left(t+z_1\lambda_1(y_1)+z_2\lambda_1(y_2),\ldots,
        t+z_1\lambda_n(y_1)+z_2\lambda_n(y_2)\right)\ne0.
 \tag{I.90}
\]
This holds in every characteristic. It improves I-36's five-pair
bound on the stated affine subclass. It does not cover a mixture
containing variable-dependent generic-rank-three factors, or
replace the affine hypothesis by arbitrary individual degree.

We prove the slightly more convenient rational statement. For
arbitrary $c_1,\ldots,c_n\in E^\times$, set
\[
 R_{k,i}=t c_i+\sum_{h=1}^{k}\frac{z_h}{a_i-y_h}.
 \tag{I.91}
\]
Every input in the promised class satisfies $f(\mathbf R_2)\ne0$.
All parameters here are algebraically independent; rational
identities are understood in their rational function field.

\paragraph{1. Pairwise dominance and differential consequences.}
Write $\alpha_i=t c_i+z/(a_i-y)$. For $i\ne j$, the determinant
\[
 \frac{c_i}{a_j-y}-\frac{c_j}{a_i-y}
 =\frac{c_i(a_i-y)-c_j(a_j-y)}{(a_i-y)(a_j-y)}
 \tag{I.92}
\]
is nonzero. If its numerator were zero, comparison of its linear
and constant coefficients would give $c_i=c_j$ and
$c_i(a_i-a_j)=0$, a contradiction. The substitution in the two
coordinates is thus an invertible linear change of $(t,z)$ over
$E(y)$. It preserves every nonzero polynomial in at most two
specified coordinates, over any fixed extension of $E$ as well.

Suppose, toward a contradiction, that $f(\mathbf R_2)=0$. Setting
the second amplitude to zero gives $f(\boldsymbol\alpha)=0$.
The coefficient of a new amplitude $\epsilon$ in
$f(\alpha_i+\epsilon/(a_i-u):1\le i\le n)=0$ is
\[
 \sum_i\frac{f_i(\boldsymbol\alpha)}{a_i-u}=0,
 \qquad f_i=\partial f/\partial X_i.
\]
The distinct simple poles force $f_i(\boldsymbol\alpha)=0$ for
every $i$: multiply by $\prod_j(a_j-u)$ and set $u=a_i$.
Since $\boldsymbol\alpha+s\mathbf c$ is obtained by replacing
$t$ by $t+s$, these identities also give
\[
 \sum_j c_j f_{ij}(\boldsymbol\alpha)=0
 \quad\hbox{for every }i,
 \qquad H(\boldsymbol\alpha)\mathbf c=0,
 \tag{I.93}
\]
where $f_{ij}$ denotes the second partial derivative. This uses
only the coefficient of $s$, so it is valid in every characteristic.
For multilinear $f$, every diagonal entry $f_{ii}$ is zero.

\paragraph{2. Irreducible factors inherit the representation promise.}
We may extend the coefficient field to its algebraic closure.
This is a proof device, not a requirement that the final sampler
evaluate over the algebraic closure. A nonzero substitution over
an extension is a nonzero substitution over its original field.
The kernel of the rational substitution $\mathbf X=\mathbf R_2$
is prime, because the target is a field. Hence an assumed nonzero
annihilator has a nonconstant irreducible factor $g$ that is also
an annihilator.

The affine-factor representation makes $f$ multilinear. For a
product of nonzero polynomials, degree in each fixed variable is
additive. Thus the nonconstant irreducible factors of a multilinear
polynomial depend on disjoint variable sets. Specialize all
variables outside the set used by $g$ so that the product of the
other factors is a nonzero scalar. Such a specialization exists
over the infinite algebraically closed field. It realizes a
nonzero scalar multiple of $g$ as a specialization of the original
matrix word. Remove the now constant factors by absorbing their
products into a neighboring factor or a boundary. Multiplication
by constant matrices preserves affine dependence and cannot
increase generic rank. The surviving factors therefore still
have generic rank at most two. A width-two or width-one word may
be padded with zero coordinates to width three.

Consequently it suffices to treat an irreducible $f$ that depends
on every one of its remaining variables, with those variables
listed in their representation order. The corresponding nodes
remain distinct and the corresponding $c_i$ remain nonzero.
The proof does not enumerate or discover this order. If at most
two variables remain, part~1 already contradicts annihilation.
Assume henceforth that $n\ge3$.

\paragraph{3. Every strict internal interval has generic rank two.}
We use an elementary polynomial rank-one factorization. If a
nonzero matrix $B$ over a multivariate polynomial ring over a
field has rank one over its fraction field, then $B=pq$ for a
polynomial column $p$ and polynomial row $q$. Indeed, take a
nonzero polynomial column, divide its entries by their greatest
common divisor, and call the primitive resulting column $p$.
Every column of $B$ is a rational multiple $h p$. In a reduced
fraction for $h$, each irreducible denominator factor would have
to divide every nonzero entry of $p$. Primitivity excludes it,
so $h$ is polynomial. This proves the asserted factorization.
In particular, when $B$ is multilinear, the polynomial factors
$p$ and $q$ depend on disjoint variable sets: the maximum entry
degrees satisfy
$\deg_{X_i}B=\deg_{X_i}p+\deg_{X_i}q\le1$.

Write the irreducible scalar as
$f=u A_1(X_1)\cdots A_n(X_n)v$. For a strict internal interval
$2\le s\le t\le n-1$, the product $B=A_s\cdots A_t$ cannot
be zero, because then $f=0$. If it had rank one, the preceding
factorization would give
\[
 f=(uA_1\cdots A_{s-1}p)(qA_{t+1}\cdots A_nv).
 \tag{I.94}
\]
The first factor depends on $X_1$ and the second on $X_n$,
because $f$ does and each of these variables appears on only
its indicated side. Both factors are therefore nonconstant,
contradicting irreducibility. Since each original factor has
rank at most two, every strict internal interval has rank
exactly two.

The pairwise-dominant map of part~1 preserves every nonzero
two-by-two minor of every such interval. To see this with the
same quantifiers as I-36, take second compounds: each
$C_2(A_i)$ has generic rank one, and Cauchy--Binet turns an
interval minor into an entry of their product. The pivot-cleared
rank-one identity (I.86) writes that entry, times univariate
nonzero pivots, as a product of univariate endpoint polynomials
and adjacent bivariate contractions. Part~1 preserves every
factor. Thus every internal interval has rank two at
$\mathbf X=\boldsymbol\alpha$, simultaneously.

\paragraph{4. Compatible constants and nonzero active endpoints.}
Work over $K=\overline E(t,z,y)$, and translate to independent
variables $Y_i=X_i-\alpha_i$. Absorb the first and last factors
into the affine endpoint row $u(X_1)=uA_1(X_1)$ and column
$v(X_n)=A_n(X_n)v$. The interior consists of positions
$2,\ldots,n-1$.

The constant row $u(\alpha_1)A_2(\alpha_2)$ is nonzero. Indeed,
the polynomial row $u(X_1)A_2(X_2)$ is nonzero, since otherwise
$f$ would be zero, and one of its nonzero bivariate entries
survives by part~1. Similarly
$A_{n-1}(\alpha_{n-1})v(\alpha_n)\ne0$. Because every internal
constant interval has rank two, multiplication by the remaining
constant interior factors is injective on these two-dimensional
active spaces. In particular both endpoint projections to the
persistent active space are nonzero.

Here is an explicit justification for the compatible gauges.
View the interior constants $C_i=A_i(\alpha_i)$ as linear maps
from the state space at cut $i$ to the state space at cut $i-1$.
At each internal cut, the image of the next constant factor is
two-dimensional and is complementary to the kernel of the
preceding constant factor, because the product of those two
factors has rank two. Use that image as the active space and
that kernel as its one-dimensional complement. Choose an
active complement to the kernel at the final cut, and transport
its basis leftwards through the constants; the full interval
rank condition makes all transported active bases invertible.
At the initial cut choose any complementary line. In these
bases every interior constant is
\[
 P=\operatorname{diag}(1,1,0),\qquad
 A_i(\alpha_i+Y_i)=P+Y_iD_i.
 \tag{I.95}
\]
The nonzero active endpoint row $a$ and column $b$ satisfy
$ab=f(\boldsymbol\alpha)=0$. Choose an active basis with
$a=e_1^{\mathsf T}$ and $b=e_2$; this is possible by taking
the second active basis vector to be $b$ and choosing the first
to have pairing one with $a$. Use this same active basis change
at every cut, keeping all copies of $P$ fixed.

The passive endpoint components can also be removed while
leaving those constants fixed. At the initial cut use
$L=\left(\begin{smallmatrix}I_2&h\\0&1\end{smallmatrix}\right)$:
$L^{-1}P=P$ and the row changes from $(a,\rho)$ to
$(a,ah+\rho)$. Choose $ah=-\rho$. At the final cut use
$R=\left(\begin{smallmatrix}I_2&0\\k^{\mathsf T}&1\end{smallmatrix}\right)$:
$PR=P$ and the column becomes
$(b,\sigma-k^{\mathsf T}b)^{\mathsf T}$. Choose
$k^{\mathsf T}b=\sigma$. These independent endpoint changes
also work when there is only one interior factor. They preserve
all matrix ranks and affine dependence. We have obtained
\[
 f(\boldsymbol\alpha+\mathbf Y)
 =(e_1^{\mathsf T}+Y_1p)
   \prod_{i=2}^{n-1}(P+Y_iD_i)
   (e_2+Y_nq).
 \tag{I.96}
\]
Here $p$ is a row and $q$ a column. All gauges are constant with
respect to the residual variables $\mathbf Y$; their dependence
on the center parameters is allowed.

Since every affine interior pencil has determinant identically
zero, the coefficient of $Y_i$ in
$\det(P+Y_iD_i)$ gives $(D_i)_{33}=0$. The first-derivative
identities from part~1 give
\[
 p_2=0,\qquad q_1=0,\qquad (D_i)_{12}=0
 \quad(2\le i\le n-1).
 \tag{I.97}
\]
Thus each active two-by-two derivative block is lower triangular,
and the passive state has no self-loop in either $P$ or $D_i$.

\paragraph{5. A tridiagonal Hessian forces every output path to vanish.}
If $|i-j|>1$, at least one constant factor $P$ separates the
two differentiated positions in (I.96). It kills the passive
coordinate. The active derivative blocks are lower triangular,
the active part of $p$ is supported on state one, and the active
part of $q$ on state two. Therefore $f_{ij}(\boldsymbol\alpha)=0$
for $|i-j|>1$. The diagonal is zero by multilinearity.

Consequently $H(\boldsymbol\alpha)$ is symmetric tridiagonal
with zero diagonal. Equation (I.93), and $c_i\ne0$ for every
$i$, force all its entries to be zero. The first row gives
$H_{12}=0$; the second row then gives $H_{23}=0$; continuing
along the path proves $H_{i,i+1}=0$ for all $i$. This argument
uses no division by two.

The adjacent entries have particularly simple formulas:
\[
 \begin{aligned}
 H_{12}&=p_3(D_2)_{32},\\
 H_{i,i+1}&=(D_i)_{13}(D_{i+1})_{32}
                 &&(2\le i\le n-2),\\
 H_{n-1,n}&=(D_{n-1})_{13}q_3.
 \end{aligned}
 \tag{I.98}
\]
All alternative active contributions are zero by (I.97).

Expand (I.96) as a sum over state paths. An active path cannot
move directly from state one to state two, because every
$(1,2)$ entry is zero. The passive state three has no self-loop.
Hence any path that first reaches state two after starting in
state one must contain a transition $1\to3$ immediately
followed by $3\to2$. Its two entries have product equal to
one of the middle expressions in (I.98), times the two residual
variables, and the product is zero. If the path starts in the
passive endpoint component, its first possible move into state
two is covered by the first expression in (I.98); if it first
returns to state one, the preceding argument applies later.
Finally, an output through the passive component of the final
column is covered by the last expression in (I.98), unless
state two was already reached, in which case an earlier zero
transition has already occurred. A path staying in state one
cannot exit through the active final column because $q_1=0$.
Thus every contributing state path contains a zero factor.
The polynomial (I.96) is identically zero over $K$, contrary
to the nonzero input. This proves $f(\mathbf R_2)\ne0$.

\paragraph{6. Conversion to ordinary centered SV and exact sampling.}
Put $A(T)=\prod_i(T-a_i)$ and $d_i=A'(a_i)\ne0$. For the
diagonally rescaled polynomial
$F(W_1,\ldots,W_n)=f(W_1/d_1,\ldots,W_n/d_n)$, the affine
singular-factor promise is unchanged. Take $c_i=d_i$ in (I.91).
The Lagrange identity
\[
 d_i\left(t+\sum_{h=1}^{2}z_h\lambda_i(y_h)\right)
 =d_i t-\sum_{h=1}^{2}\frac{z_h A(y_h)}{a_i-y_h}
 \tag{I.99}
\]
shows that the ordinary centered map is the rational map just
proved, after the invertible change of amplitudes
$z_h\mapsto-z_hA(y_h)$ in the rational parameter field.
Therefore its polynomial substitution is nonzero, proving (I.90).

For total degree at most $D$ (here $D\le n$), that substitution
has degree at most $D$ in $t$, at most $(n-1)D$ in each selector,
and total degree at most $D$ in the two amplitudes. The Newton
simplex interpolation argument already proved in I-21 yields
a hitting set of at most
\[
 (D+1)\bigl((n-1)D+1\bigr)^2\binom{D+2}{2}.
 \tag{I.100}
\]
The working field must supply at least
$\max\{n,D+1,(n-1)D+1\}$ distinct elements. A proper working
extension uses the same extension-evaluation convention as
I-21 and I-36. The proof holds in all characteristics, and
$D=0$ requires just one evaluation. The selected order and
all coefficient-dependent factorizations occur only in the
proof, not in the sampler.\hfill$\square$

\paragraph{Provenance and the remaining obstruction.}
This deduction uses the pairwise dominance and rank-one compound
factorization proved in I-36. The first-order pole-extraction
observation was supplied independently in this run by the purity
agent while investigating the Hu--van Melkebeek--Morgan
vanishing-ideal framework; the full argument above proves the
needed arbitrary-slope and arbitrary-characteristic version
directly. No historical priority claim is made.

The new closing step is specific to affine singular pencils:
after rank preservation their passive constant entry and passive
linear coefficient both vanish. A variable-dependent invertible
factor can propagate the passive state, so nonadjacent Hessian
entries need not vanish and (I.93) need not force all adjacent
transition products to vanish. Higher-degree singular factors
also need not have a zero passive entry at every positive degree.
Thus the mixed rank-two/rank-three centered-two-pair gate remains
open, as does an unrestricted polynomial-in-width construction.
